% Guía del instructor, edición 3.0 (banco completo).

\section{Resumen operativo}

\begin{resumenejecutivo}
\begin{itemize}
  \item \textbf{Dos laboratorios de solo recepción}: Lab.\ 01 (Windows, FM, 210 min) y Lab.\ 02 (Ubuntu 24.04, ADS-B/Mode S, 225 min).
  \item \textbf{Cuatro documentos}: el \emph{manual} (procedimiento), el \emph{Anexo A} (fundamentos y diagnóstico), el \emph{Anexo B} (cuestionarios y hojas de registro, que el estudiante imprime) y esta guía.
  \item \textbf{Banco completo}: HackRF One, ANT500 para FM, antena de 1090 MHz con base magnética, filtro pasivo de 1090 MHz, cable de extensión SMA, adaptadores SMA, adaptador N -- SMA, SIGLENT SSA3032X Plus y atenuador SMA de 20 dB. La lista con enlaces de compra está en el README del repositorio.
  \item \textbf{Antes de cada curso}: instalar el software y hacer el ensayo previo de la sección~\ref{sec:i-ensayo}.
\end{itemize}
\end{resumenejecutivo}

\section{Requisitos del curso y del estudiante}

\begin{itemize}
  \item \textbf{Conceptuales:} frecuencia y longitud de onda, dB y dBm, ancho de banda, muestreo, AM/FM, radar primario y secundario, binario y hexadecimal. Los cuestionarios diagnósticos detectan carencias: si más de un tercio del grupo falla las preguntas 1--3 del diagnóstico del Lab.\ 01, dedique 10 minutos a dB y longitud de onda durante la orientación.
  \item \textbf{Prácticos:} copiar órdenes con exactitud, abrir una terminal y capturar pantalla. No se asume experiencia con SDR ni con Linux.
  \item \textbf{Organización:} grupos de 2 por banco. Con un solo SIGLENT, rote la etapa 5 del Lab.\ 01 con la etapa 6, y la etapa 8 del Lab.\ 02 con la etapa 7.
\end{itemize}

\section{Tiempo del instructor}

\begin{longtable}{@{}>{\raggedright\arraybackslash}p{0.55\linewidth}>{\raggedleft\arraybackslash}p{0.15\linewidth}>{\raggedright\arraybackslash}p{0.22\linewidth}@{}}
\toprule
\textbf{Actividad} & \textbf{Tiempo} & \textbf{Cuándo} \\
\midrule
\endhead
Instalación de software por puesto & 30--45 min & Una vez por curso \\
Ensayo previo del Lab.\ 01 & 45 min & Antes del curso \\
Ensayo previo del Lab.\ 02 (vivo y grabación) & 60 min & Antes del curso, a la hora de la clase \\
Comprobación rápida por puesto & 10 min & El día de la sesión \\
\bottomrule
\end{longtable}

\section{El banco y su preparación}

\begin{checklist}
  \item HackRF One con su cable USB de datos; los tres SMA limpios.
  \item ANT500 extensible sin holguras; antena de 1090 MHz con su cable de 3 m y base magnética.
  \item Filtro de 1090 MHz: pasivo (no requiere alimentación); géneros de sus conectores compatibles con la cadena, con los adaptadores SMA del paquete del HackRF.
  \item Adaptador N macho -- SMA hembra en cada analizador.
  \item Atenuador SMA de 20 dB guardado junto al banco (solo para emplazamientos muy cercanos a radares; véase la sección~\ref{sec:i-siglent}).
  \item SIGLENT: \gui{Preset} correcto, TG apagado y fecha de calibración anotada.
  \item Una ventana con cielo despejado por aula para la antena de 1090 MHz, con un recorrido de cable seguro.
  \item Grabación de referencia copiada en la carpeta personal de cada usuario de Ubuntu.
\end{checklist}

\section{Software en Windows}

\begin{longtable}{@{}>{\raggedright\arraybackslash}p{0.30\linewidth}>{\raggedright\arraybackslash}p{0.62\linewidth}@{}}
\toprule
\textbf{Componente} & \textbf{Instalación y comprobación} \\
\midrule
\endhead
SDRangel 7.27.2 & Instalador oficial \texttt{sdrangel-7.27.2-win64.exe} (\url{https://github.com/f4exb/sdrangel/releases}), SHA-256 \texttt{96e4d6c83ad554a7 8173310f78302162 fd52463d6a890da5 8211021dabf74dc2}. Incluye el runtime de Visual C++. \\
Herramientas HackRF & radioconda (\url{https://github.com/radioconda/radioconda-installer}), recomendado por la documentación de HackRF; aporta \texttt{hackrf\_info}. \\
Controlador & «HackRF One» sin advertencias en el Administrador de dispositivos. Zadig solo si falta (procedimiento en el Anexo A). \\
Comprobación & \tecla{Ctrl}+\tecla{R} en SDRangel muestra el HackRF con su número de serie. \\
\bottomrule
\end{longtable}

\section{Software en Ubuntu 24.04}

\begin{enumerate}
  \item \texttt{sudo apt update \&\& sudo apt install hackrf} (paquete 2023.01.1-9build1, \emph{universe}); confirme la regla \texttt{/usr/lib/udev/rules.d/60-libhackrf0.rules}.
  \item \texttt{hackrf\_info} sin \texttt{sudo} debe encontrar el HackRF; si no, siga el Anexo A (permisos).
  \item \texttt{sudo snap install sdrangel} y \texttt{sudo snap connect sdrangel:raw-usb}. No use el \texttt{.deb} de 24.04 de la versión 7.27.2: depende de \texttt{libuhd4.10.0}, que no está en los repositorios estándar.
  \item Copie \texttt{datos\_adsb/} a la carpeta personal de cada usuario y ábrala una vez con \textbf{FileInput} para confirmar el acceso desde el snap.
\end{enumerate}

\section{El SIGLENT y los emplazamientos cercanos a radares}
\label{sec:i-siglent}

\begin{itemize}
  \item Límites (hoja de datos DS0703P): RF INPUT N hembra de 50 $\Omega$; potencia media máxima de 30 dBm durante 3 min con $f_c \ge 10$ MHz, atenuación $> 20$ dB y preamplificador apagado; daño a 33 dBm; DC máximo $\pm 50$ V. Con antenas pasivas el riesgo real es la descarga electrostática.
  \item Rótulos de las teclas, \gui{N dB BW}, \gui{Demod} y tiempo de barrido mínimo (900 $\mu$s): manual de usuario UM0703P-E01A. Las funciones \emph{Occupied BW} y \emph{Spectrum Monitor} pertenecen al kit opcional AMK.
  \item \textbf{Radares cercanos.} El HackRF no tiene preselección. A menos de unos pocos kilómetros de un SSR o un DME, el pico recibido puede acercarse a $-5$ dBm (estimación en el Anexo A). En ese caso prohíba el \gui{RF amp}, no oriente la antena hacia el radar y, si el emplazamiento es muy cercano, intercale el atenuador de 20 dB entre el filtro y el HackRF, informando al grupo de que perderá sensibilidad.
\end{itemize}

\section{HackRF: firmware}

Anote con \texttt{hackrf\_info} la versión de firmware y la del software. Herramientas oficiales actuales: v2026.01.3; Ubuntu 24.04 empaqueta 2023.01.1. Si algún programa advierte de incompatibilidad, actualice el firmware siguiendo la documentación oficial de HackRF, fuera de clase: una actualización interrumpida deja la placa en modo DFU.

\section{Versiones de referencia}

\begin{longtable}{@{}>{\raggedright\arraybackslash}p{0.30\linewidth}>{\raggedright\arraybackslash}p{0.62\linewidth}@{}}
\toprule
\textbf{Componente} & \textbf{Versión} \\
\midrule
\endhead
SDRangel & 7.27.2 (Windows: instalador oficial; Ubuntu: snap) \\
Herramientas HackRF & 2026.01.3 (radioconda) / 2023.01.1 (Ubuntu) \\
Sistemas & Windows 10/11; Ubuntu 24.04 LTS nativo \\
Verificador de la grabación & Python 3.12 o superior con \texttt{numpy} \\
\bottomrule
\end{longtable}

\section{Ensayo previo}
\label{sec:i-ensayo}

Siga el manual \emph{tal como está escrito}, en el aula y a la hora de la clase. Anote los resultados en la tabla y corrija el manual si su versión de software muestra algo distinto.

\begin{longtable}{@{}>{\raggedright\arraybackslash}p{0.50\linewidth}>{\raggedright\arraybackslash}p{0.42\linewidth}@{}}
\caption{Registro del ensayo previo.}\label{tab:i-ensayo}\\
\toprule
\textbf{Comprobación} & \textbf{Resultado y fecha} \\
\midrule
\endfirsthead
\toprule
\textbf{Comprobación} & \textbf{Resultado y fecha} \\
\midrule
\endhead
Lab.\ 01, etapas 1--4: audio inteligible de al menos 3 emisoras & \\[2mm]
Lab.\ 01, etapa 5: emisora medida con marcador y $-26$ dB & \\[2mm]
Lab.\ 02, etapas 1--3: tramas válidas en 10 min & \\[2mm]
Lab.\ 02, etapa 6: aeronaves con posición en el mapa & \\[2mm]
Lab.\ 02, etapa 8: tiempo de barrido mínimo en span cero; ráfagas capturadas & \\[2mm]
Lab.\ 02, etapa 9: grabación con valores de referencia & \\[2mm]
Distancia del aula al SSR o DME más cercano & \\[2mm]
\bottomrule
\end{longtable}

\section{Criterios de éxito}

\begin{longtable}{@{}>{\raggedright\arraybackslash}p{0.12\linewidth}>{\raggedright\arraybackslash}p{0.41\linewidth}>{\raggedright\arraybackslash}p{0.39\linewidth}@{}}
\toprule
\textbf{Lab.} & \textbf{Mínimo para aprobar} & \textbf{Excelencia} \\
\midrule
\endhead
01 & HackRF identificado; receptor según la etapa 3; emisora localizada; demodulador con desplazamiento y de-énfasis correctos; audio; analizador con marcador; hoja W8 razonada. & Cuatro experimentos con predicción y explicación física; ancho a $-26$ dB coherente con la regla de Carson; límites de la comparación de niveles bien argumentados. \\
02 & Ubuntu y HackRF verificados; tramas válidas; trama decodificada a mano; tabla con el origen de cada campo; mapa; grabación de referencia con la lección CPR explicada. & Experimento A--B--A concluyente o declarado no concluyente con razones; span cero con duración medida; comparación vivo/grabación. \\
\bottomrule
\end{longtable}

\begin{importante}[Lo que nunca debe premiarse]
Tablas completas con datos que el grupo no pudo observar, capturas de servicios web presentadas como propias o niveles en dBm comparados sin matices con dB de SDRangel.
\end{importante}

\section{Conducción de la sesión}

\begin{longtable}{@{}>{\raggedright\arraybackslash}p{0.22\linewidth}>{\raggedright\arraybackslash}p{0.70\linewidth}@{}}
\toprule
\textbf{Momento} & \textbf{Qué vigilar} \\
\midrule
\endhead
Lab.\ 01, etapa 1 & Antena en CLKIN o CLKOUT: el receptor funciona, pero no se ve ninguna emisora. \\
Lab.\ 01, etapa 3 & Grupos que no cierran radioconda Prompt antes de SDRangel. \\
Lab.\ 01, etapa 4 & Desplazamiento con signo positivo; de-énfasis en 50 $\mu$s; confusión del pico DC con una emisora. Pregunte: «¿qué demuestra ese pico?». \\
Lab.\ 01, etapa 6 & Grupos que cambian dos cosas a la vez; insista en una variable por prueba. \\
Lab.\ 02, etapa 3 & Si a los 10 minutos no hay tramas en ningún puesto, el cielo puede estar vacío: adelante la etapa 9 y vuelva después. \\
Lab.\ 02, etapa 6 & Capturas con coordenadas exactas: recuerde la regla de privacidad. \\
Lab.\ 02, etapa 9 & No anticipe la respuesta: deje que el grupo vea la aeronave en un lugar imposible antes de fijar la referencia. Es el momento de más aprendizaje del laboratorio. \\
\bottomrule
\end{longtable}

\section{Problemas frecuentes}

\begin{longtable}{@{}>{\raggedright\arraybackslash}p{0.30\linewidth}>{\raggedright\arraybackslash}p{0.62\linewidth}@{}}
\toprule
\textbf{Problema} & \textbf{Recuperación} \\
\midrule
\endhead
Ninguna trama en 10 min & Revisar la cadena (antena, filtro, adaptadores); antena en la ventana; LNA 24--32. Si todo está bien, adelantar la etapa 9. No bajar el umbral de 5 dB para «conseguir» filas. \\
Aeronave de la grabación en un lugar imposible & Esperado con la posición del aula. Desactivar \gui{Auto-update position} y fijar 37,0 / 14,0. \\
FileInput no abre el archivo (Ubuntu) & Archivo fuera de la carpeta personal o alterado: copiarlo a \texttt{\textasciitilde/} y comprobar el SHA-256. \\
Hay emisora pero no audio & Signo del desplazamiento, mute, squelch o dispositivo de audio. \\
HackRF «ocupado» & Otro programa abierto. \\
Saturación con RF amp & Apagarlo; aproveche para mostrar los espurios. \\
SDRangel lento & Mapa 2D; no activar \emph{Display All Flight Paths}. \\
\bottomrule
\end{longtable}

\section{La grabación de referencia}

\begin{enumerate}
  \item Integridad: \texttt{sha256sum} debe dar {\ttfamily\small\nolinkurl{d201250f486e5ff14133bc57664c8e5be6c42be3bb989529f29789674ee5f25d}}.
  \item Opcional: \texttt{python3 herramientas/adsb\_iq.py verificar \ldots} informa 430 preámbulos candidatos, 242 tramas válidas y 1 dirección ICAO.
  \item En SDRangel 7.27.2 (Windows), FileInput con bucle y el ADS-B Demodulator con \gui{SR} 2M produjeron la fila \texttt{4d2023}/AMC421: 20\,025--24\,275 ft, 371--390 kt y posiciones 36,96--37,10°N, 13,79--13,86°E con la referencia 37,0/14,0. Con la posición tomada de Windows (Bogotá), la aeronave apareció en Colombia.
  \item La casilla \gui{Auto-update position} viene activada por defecto y sobrescribe cada segundo la posición fijada: desactívela.
\end{enumerate}

\section{Privacidad y uso de datos}

\begin{itemize}
  \item Las capturas publicadas no deben mostrar las coordenadas del receptor.
  \item \gui{Feed} se mantiene desactivado salvo autorización institucional.
  \item La grabación contiene datos públicos radiodifundidos en 2013 por una aeronave real; se usa con fines educativos bajo la licencia BSD-3 de dump1090.
\end{itemize}

\section{Claves de los cuestionarios}
\label{sec:i-claves}

\subsection{Laboratorio 01 · diagnóstico}
\CuestionarioDiagnosticoW
\subsection{Laboratorio 01 · final}
\CuestionarioFinalW
\subsection{Laboratorio 02 · diagnóstico}
\CuestionarioDiagnosticoU
\subsection{Laboratorio 02 · final}
\CuestionarioFinalU

\section{Respuestas de referencia de las hojas}

\begin{longtable}{@{}>{\raggedright\arraybackslash}p{0.17\linewidth}>{\raggedright\arraybackslash}p{0.75\linewidth}@{}}
\toprule
\textbf{Hoja} & \textbf{Respuesta esperada} \\
\midrule
\endhead
W2 & Varias emisoras de alturas distintas: potencia, distancia, obstáculos y respuesta de la antena difieren. \\
W3 & Nivel 2 (audio inteligible); nivel 5 solo si se decodificó RDS. \\
W4 & Ancho a $-26$ dB del orden de 150--300 kHz; nivel en dBm dependiente de la antena y del adaptador. \\
W5 & 50 kHz: distorsión (parte de la desviación fuera del filtro); 100 kHz: muy degradado; 200 kHz: emisora vecina o ruido. \\
W6 & RF BW 80 kHz: distorsión por recorte; 250 kHz: más ruido y posible interferencia adyacente. AF BW 3 kHz: audio «telefónico» pero inteligible. RF BW actúa antes de demodular y AF BW después. \\
W7 & Baja: emisora apenas sobre el ruido. Moderada: mejor compromiso. Alta y excesiva: piso más alto, espurios y distorsión. \\
W8 & Espectro: ambos sí. dBm con exactitud especificada: SIGLENT sí; SDRangel no. Audio: SDRangel sí; SIGLENT también con \gui{Demod} (acepte «sí» o «parcial» si se justifica). Cascada: SDRangel sí; SIGLENT parcial (\emph{Spectrum Monitor} del kit AMK). Ancho con marcadores: SIGLENT sí; SDRangel parcial. 20 MHz: el SIGLENT los recorre en cada barrido, no simultáneamente; el HackRF los vería a la vez a 20 MSa/s, no con 2,4 MSa/s. Antena y cable: ambos sí. \\
W9 & Niveles 1 y 2 (y 5 con RDS); Carson con las medidas de W4; referencias distintas (escala digital frente a dBm), ganancias y RBW distintas, sin calibración. \\
U4 & Un pico es nivel 1: puede ser ruido u otro emisor; solo la paridad coherente con el formato valida los bits. \\
U5 & \texttt{8D} = \texttt{10001101}: DF 17, CA 5. ICAO \texttt{4D2023}. ME \texttt{20} = \texttt{00100000}: TC 4, categoría 0. Grupos: \texttt{000001}, \texttt{001101}, \texttt{000011}, \texttt{110100}, \texttt{110010}, \texttt{110001}, \texttt{100000}, \texttt{100000} = 1, 13, 3, 52, 50, 49, 32, 32 = «AMC421». Es un mensaje de identificación: no contiene posición, y una sola trama de posición tampoco bastaría sin un par o una referencia. \\
U6 & ICAO, indicativo, altitud y velocidad: T. Lat/Lon, distancia y acimut: D. Modelo y matrícula: E. La identificación y la altitud llegan en mensajes distintos, con ritmos distintos. \\
U7 & La posición (decodificación local), la distancia y el alcance dependen de \gui{My Position}; los mosaicos del mapa, de Internet. \\
U8 & Se valora el razonamiento: si A' difiere de A más que B de A, el efecto de la ganancia no está demostrado. El umbral más bajo (C) produce más tramas y más inválidas; sin filtro (D) puede bajar la calidad cerca de emisoras intensas. \\
U9 & El analizador solo ve energía en función del tiempo (nivel 1); no demodula los bits ni comprueba la paridad. \\
U10 & La trama transmite la posición dentro de una zona de unos 6°; la decodificación local eligió la zona más cercana a la referencia del aula. Nada en la trama permite detectar el error. \\
U11 & Niveles 4--6 con evidencia propia; mejoras justificadas con U8 (ubicación antes que ganancia); modelo, matrícula y mosaicos no proceden de la radio. \\
\bottomrule
\end{longtable}

\section{Rúbrica}

La rúbrica para el estudiante está en el Anexo B (sección «Cómo se evalúa»). Para cada criterio, asigne el máximo cuando la evidencia esté registrada, justificada y etiquetada con su nivel; la mitad cuando esté registrada sin justificar; cero cuando falte o no pueda haberse observado.

\section{Limitaciones científicas}

\begin{itemize}
  \item Ningún nivel medido es trazable: antenas y cables no están calibrados; el HackRF no es un instrumento de medida.
  \item La recepción a 1090 MHz depende del lugar, la hora y el tráfico: dos sesiones con el mismo equipo pueden dar resultados muy distintos.
  \item La grabación de referencia tiene el tiempo comprimido y se repite en bucle: no sirve para estudiar tasas, alcance ni seguimiento.
  \item El span cero muestra envolventes, no bits.
  \item El verificador del repositorio es deliberadamente simple: encuentra menos tramas que SDRangel (126 comunes; SDRangel registró 17 DF17 adicionales; el verificador, 1 DF21 que SDRangel no registró).
\end{itemize}
